\documentclass[12pt,a4paper]{article}
\usepackage[margin=2.3cm]{geometry}
\usepackage{xcolor}
\usepackage[T1]{fontenc}
\usepackage[utf8]{inputenc}
\usepackage[brazil]{babel}
\usepackage{lmodern}
\usepackage{amsmath,amssymb,booktabs,array}
\usepackage{tikz}
\usetikzlibrary{automata,positioning,arrows.meta}
\usepackage{listings}
\definecolor{azul}{RGB}{20,49,78}
\definecolor{verde}{RGB}{0,121,111}
\definecolor{cinza}{RGB}{244,247,250}
\lstset{language=Python,basicstyle=\ttfamily\scriptsize,
 keywordstyle=\color{azul}\bfseries,commentstyle=\color{verde},
 stringstyle=\color{verde},breaklines=true,showstringspaces=false,
 frame=single,rulecolor=\color{azul!20},backgroundcolor=\color{cinza},
 columns=fullflexible,keepspaces=true,tabsize=4}
\newcommand{\q}[1]{q_{\mathrm{#1}}}
\newcommand{\MiniL}{\textsc{MiniL}}
\newcommand{\grupo}{Gabriel Holanda, Gabriel Torres, Victor Milito,\\Paulo Júnior, Vinicius Kaua e Pablo Café}
\tikzset{every state/.style={draw=azul,thick,minimum size=8mm,font=\small},
 every edge/.style={draw=azul,>=Stealth},every edge quotes/.style={font=\scriptsize}}

\usepackage{microtype}
\usepackage{needspace}
\usepackage[colorlinks=true,linkcolor=azul,urlcolor=verde]{hyperref}
\usepackage{fancyhdr}
\pagestyle{fancy}\fancyhf{}\fancyhead[L]{\small\MiniL{} | Reconhecedor léxico}
\fancyhead[R]{\small LFAF}\fancyfoot[C]{\thepage}
\setlength{\headheight}{15pt}
\setlength{\parindent}{0pt}\setlength{\parskip}{6pt}
\setlength{\emergencystretch}{2em}
\setcounter{tocdepth}{1}
\lstset{basicstyle=\ttfamily\footnotesize}
\title{\color{azul}\textbf{O Compilador Mínimo: relatório técnico}\\[4mm]
\large Modelando um reconhecedor léxico para a mini-linguagem \MiniL}
\author{\normalsize\grupo}
\date{Linguagens Formais e Autômatos Finitos\\Revisão 1.1.0 --- 13 de novembro de 2026}
\begin{document}
\maketitle
\begin{abstract}
Este trabalho especifica e implementa um analisador léxico para uma mini-linguagem
com palavras-chave, identificadores, inteiros, decimais, operadores e separadores.
Definimos um AFD completo com nove estados, apresentamos sua tabela e seus
diagramas vetoriais, e aplicamos o refinamento de partições de Moore, obtendo
oito estados no reconhecedor de lexemas. A implementação em Python utiliza a
mesma função de transição e a regra do maior prefixo aceito. O conjunto de
testes verifica tokens, erros, posições, alcançabilidade e equivalência exata
entre os autômatos, por exploração do produto. O projeto cobre a etapa léxica;
análise sintática, semântica e geração de código ficam fora do escopo.
\end{abstract}
\textbf{Palavras-chave:} AFD; análise léxica; expressões regulares; minimização; LaTeX.
\tableofcontents
\newpage
\section{Problema, objetivo e decisões de projeto}
O scanner recebe um texto e o transforma em uma sequência de tokens. Por exemplo,
\texttt{let total = 12.50;} produz uma palavra-chave, um identificador, um operador,
um literal decimal e um separador. Cada token armazena categoria, lexema, linha e
coluna. Lexema é o texto concreto; categoria é sua classificação.

A aplicação escolhida é \emph{Analisador Léxico para uma Mini-Linguagem}, prevista
na proposta da disciplina. Embora o título da proposta mencione compilador,
este projeto implementa apenas seu reconhecimento léxico. Uma sequência de
tokens válida não garante que exista um programa sintaticamente correto.

\subsection{Contrato da mini-linguagem}
\begin{itemize}
\item A linguagem diferencia maiúsculas e minúsculas: \texttt{if} é reservada;
\texttt{If} é um identificador.
\item Identificadores usam letras ASCII ou sublinhado no início; depois podem
conter também dígitos. Letras acentuadas não pertencem a essa classe.
\item Números são não sinalizados. Um decimal exige dígitos dos dois lados do
ponto. Zeros à esquerda são permitidos. Expoentes e vírgula decimal não existem.
\item Os sinais de \texttt{-3} geram dois tokens: operador e inteiro.
\item Operadores e separadores têm um caractere. \texttt{==} gera dois operadores
\texttt{=}; não existe operador de igualdade composto neste contrato.
\item Espaços, tabulações, retorno de carro e quebras de linha são reconhecidos
e descartados. A quebra de linha \texttt{\textbackslash n} avança a linha; os demais
caracteres avançam a coluna em uma unidade, inclusive tabulação.
\item A regra é o maior prefixo aceito. Assim, \texttt{12abc} gera inteiro
\texttt{12} e identificador \texttt{abc}. Não se impõe fronteira adicional.
\item Strings e caracteres não previstos causam erro léxico. Comentários não são
interpretados: barras consecutivas continuam sendo operadores, e as palavras
seguintes continuam sendo tokens comuns.
\end{itemize}

\section{Especificação por expressões regulares}
Usamos classes disjuntas de caracteres. O alfabeto finito $\Sigma$ corresponde
aos valores de caracteres representáveis por Python, de U+0000 a U+10FFFF.
Somente as classes explicitamente permitidas podem integrar um lexema aceito.
\begin{center}
\begin{tabular}{cl}
\toprule Classe & Definição\\\midrule
$L$ & \texttt{A--Z}, \texttt{a--z} e \texttt{\_}\\
$D$ & \texttt{0--9}\\
$P$ & apenas o ponto \texttt{.} (coluna ``$.$'' nas tabelas)\\
$O$ & \texttt{+}, \texttt{-}, \texttt{*}, \texttt{/}, \texttt{=}\\
$S$ & \texttt{(}, \texttt{)}, \texttt{;}\\
$W$ & espaço, \texttt{\textbackslash t}, \texttt{\textbackslash r}, \texttt{\textbackslash n}\\
$X$ & $\Sigma\setminus(L\cup D\cup P\cup O\cup S\cup W)$\\\bottomrule
\end{tabular}
\end{center}
Em regex, $*$ significa zero ou mais ocorrências e $+$ significa uma ou mais.
Na expressão $O$, o símbolo literal \texttt{+} é um membro da classe, não o
operador de repetição.
\begin{align*}
\mathrm{ID}&=L(L\mid D)^* &\mathrm{INT}&=D^+\\
\mathrm{FLOAT}&=D^+\,\texttt{.}\,D^+ &\mathrm{OP}&=O\\
\mathrm{SEP}&=S &\mathrm{WS}&=W^+.
\end{align*}
A linguagem de \emph{um lexema} é
\[
\mathcal{L}_{lex}=\mathrm{ID}\cup\mathrm{INT}\cup\mathrm{FLOAT}
\cup\mathrm{OP}\cup\mathrm{SEP}\cup\mathrm{WS}.
\]
Palavras-chave pertencem à forma regular de ID. Após consumir o maior
identificador, consultamos a tabela
$\{\texttt{let},\texttt{if},\texttt{else},\texttt{print}\}$.
Por isso, \texttt{ifx} é um único ID, e não \texttt{if} seguido de \texttt{x}.

\section{Definição formal do AFD}
De acordo com a definição clássica de AFD~\cite{sipser}, seja
\[
M=(Q,\Sigma,\delta,\q{0},F),
\]
com
\begin{align*}
Q&=\{\q{0},\q{I},\q{N},\q{Dot},\q{F},\q{Op},\q{Sep},\q{W},\q{Err}\},\\
F&=\{\q{I},\q{N},\q{F},\q{Op},\q{Sep},\q{W}\},\\
\delta&:Q\times\Sigma\longrightarrow Q.
\end{align*}
$\q{0}$ é inicial; $\q{Err}$ é o estado de erro absorvente. Os nomes dos estados
são rótulos: em particular, $\q{F}$ é o estado de decimal e $F$ é o conjunto
de estados finais. Não são o mesmo objeto.

\Needspace{18\baselineskip}
\subsection{Função de transição completa}
Cada coluna representa todos os caracteres da classe correspondente. Todas
as células estão definidas, portanto o AFD é completo e determinístico.
\begin{center}\small
\begin{tabular}{c|ccccccc}
\toprule Estado & $L$ & $D$ & $.$ & $O$ & $S$ & $W$ & $X$\\\midrule
$\q{0}$ & $\q{I}$ & $\q{N}$ & $\q{Err}$ & $\q{Op}$ & $\q{Sep}$ & $\q{W}$ & $\q{Err}$\\
$\q{I}$ & $\q{I}$ & $\q{I}$ & $\q{Err}$ & $\q{Err}$ & $\q{Err}$ & $\q{Err}$ & $\q{Err}$\\
$\q{N}$ & $\q{Err}$ & $\q{N}$ & $\q{Dot}$ & $\q{Err}$ & $\q{Err}$ & $\q{Err}$ & $\q{Err}$\\
$\q{Dot}$ & $\q{Err}$ & $\q{F}$ & $\q{Err}$ & $\q{Err}$ & $\q{Err}$ & $\q{Err}$ & $\q{Err}$\\
$\q{F}$ & $\q{Err}$ & $\q{F}$ & $\q{Err}$ & $\q{Err}$ & $\q{Err}$ & $\q{Err}$ & $\q{Err}$\\
$\q{Op}$ & $\q{Err}$ & $\q{Err}$ & $\q{Err}$ & $\q{Err}$ & $\q{Err}$ & $\q{Err}$ & $\q{Err}$\\
$\q{Sep}$ & $\q{Err}$ & $\q{Err}$ & $\q{Err}$ & $\q{Err}$ & $\q{Err}$ & $\q{Err}$ & $\q{Err}$\\
$\q{W}$ & $\q{Err}$ & $\q{Err}$ & $\q{Err}$ & $\q{Err}$ & $\q{Err}$ & $\q{W}$ & $\q{Err}$\\
$\q{Err}$ & $\q{Err}$ & $\q{Err}$ & $\q{Err}$ & $\q{Err}$ & $\q{Err}$ & $\q{Err}$ & $\q{Err}$\\\bottomrule
\end{tabular}\end{center}
A extensão a palavras satisfaz
$\delta^*(q,\varepsilon)=q$ e
$\delta^*(q,wa)=\delta(\delta^*(q,w),a)$.
Uma palavra $w$ é aceita se e somente se $\delta^*(\q{0},w)\in F$.
Como $\q{0}\notin F$, o lexema vazio não é aceito.

\subsection{Diagrama original em TikZ}
\begin{center}\begin{tikzpicture}[->,auto,node distance=17mm and 20mm,scale=.92,transform shape]
\node[state,initial] (s) {$\q{0}$};
\node[state,accepting,above right=of s] (i) {$\q{I}$};
\node[state,accepting,right=of s] (n) {$\q{N}$};
\node[state,right=of n] (d) {$\q{Dot}$};
\node[state,accepting,right=of d] (f) {$\q{F}$};
\node[state,accepting,below right=of s] (o) {$\q{Op}$};
\node[state,accepting,right=of o] (p) {$\q{Sep}$};
\node[state,accepting,below=of s] (w) {$\q{W}$};
\node[state,right=of p] (e) {$\q{Err}$};
\path (s) edge node {$L$} (i)
 (i) edge[loop above] node {$L,D$} (i)
 (s) edge node {$D$} (n)
 (n) edge[loop above] node {$D$} (n)
 (n) edge node {$.$} (d)
 (d) edge node {$D$} (f)
 (f) edge[loop above] node {$D$} (f)
 (s) edge node[left] {$O$} (o)
 (s) edge[bend left=12] node[above,pos=.75] {$S$} (p)
 (s) edge node[left] {$W$} (w)
 (w) edge[loop left] node {$W$} (w)
 (e) edge[loop right] node {$\Sigma$} (e);
\end{tikzpicture}
\end{center}
\textbf{Convenção do desenho:} toda transição não desenhada vai para
$\q{Err}$. Seu laço representa todos os caracteres de $\Sigma$. Essa convenção
evita cruzamentos; a tabela acima explicita todas as arestas omitidas.
Círculo duplo indica aceitação e a seta sem origem indica o estado inicial.

\subsection{Por que o AFD reconhece a especificação?}
O invariante de cada ramo descreve os prefixos que chegam ao seu estado:
\begin{center}
\begin{tabular}{cl}
\toprule Estado & Forma dos prefixos que o alcançam\\\midrule
$\q{0}$ & $\varepsilon$\\
$\q{I}$ & $L(L\mid D)^*$\\
$\q{N}$ & $D^+$\\
$\q{Dot}$ & $D^+\texttt{.}$\\
$\q{F}$ & $D^+\texttt{.}D^+$\\
$\q{Op}$, $\q{Sep}$ & um caractere de $O$, um de $S$, respectivamente\\
$\q{W}$ & $W^+$\\
$\q{Err}$ & prefixo que já saiu de todas as formas permitidas\\\bottomrule
\end{tabular}\end{center}
Essas propriedades seguem por indução no comprimento do prefixo, usando a
tabela. Os estados finais correspondem exatamente às seis expressões
regulares. Portanto $\mathcal{L}(M)=\mathcal{L}_{lex}$.

\section{Minimização: nove estados para oito}
Minimizamos o AFD como \textbf{reconhecedor booleano de um lexema}. Usamos
refinamento de partições de Moore: dois estados permanecem juntos quando têm
a mesma condição de aceitação e, em cada classe de entrada, chegam ao mesmo
bloco da partição anterior.

\subsection{Passo 1: remover estados inacessíveis}
Todos são alcançáveis a partir de $\q{0}$. Testemunhas:
\[
\begin{array}{c|ccccccccc}
q&\q{0}&\q{I}&\q{N}&\q{Dot}&\q{F}&\q{Op}&\q{Sep}&\q{W}&\q{Err}\\
w&\varepsilon&\texttt{a}&\texttt{0}&\texttt{0.}&\texttt{0.0}&\texttt{+}&\texttt{;}&\text{espaço}&\texttt{@}
\end{array}
\]
Nenhum estado é removido nessa etapa.

\subsection{Passo 2: partição inicial e assinaturas}
Sejam $A=F$ e $B=Q\setminus F$. Então
\[
\mathcal{P}_0=\{A,B\}=
\big\{\{\q{I},\q{N},\q{F},\q{Op},\q{Sep},\q{W}\},
\{\q{0},\q{Dot},\q{Err}\}\big\}.
\]
Na ordem $(L,D,.,O,S,W,X)$, a assinatura de $q$ é o vetor de blocos
que contêm seus destinos:
\begin{center}\small
\begin{tabular}{ccl}
\toprule Bloco original & Estado & Assinatura em $\mathcal{P}_0$\\\midrule
$A$ & $\q{I}$ & $(A,A,B,B,B,B,B)$\\
$A$ & $\q{N}$ & $(B,A,B,B,B,B,B)$\\
$A$ & $\q{F}$ & $(B,A,B,B,B,B,B)$\\
$A$ & $\q{Op},\q{Sep}$ & $(B,B,B,B,B,B,B)$\\
$A$ & $\q{W}$ & $(B,B,B,B,B,A,B)$\\
$B$ & $\q{0}$ & $(A,A,B,A,A,A,B)$\\
$B$ & $\q{Dot}$ & $(B,A,B,B,B,B,B)$\\
$B$ & $\q{Err}$ & $(B,B,B,B,B,B,B)$\\\bottomrule
\end{tabular}\end{center}
Só se comparam assinaturas dentro do mesmo bloco. Assim, $\q{Dot}$ e $\q{N}$
não se unem, pois um é não final e o outro é final. A primeira rodada produz
\[
\mathcal{P}_1=\big\{\{\q{I}\},\{\q{N},\q{F}\},\{\q{Op},\q{Sep}\},
\{\q{W}\},\{\q{0}\},\{\q{Dot}\},\{\q{Err}\}\big\}.
\]
\textbf{Atenção:} $\q{N}$ e $\q{F}$ ainda não se separam em $\mathcal{P}_1$:
ambos iam para o bloco não final ao ler ponto na partição inicial.

\subsection{Passo 3: segunda rodada e estabilização}
Em $\mathcal{P}_1$, $\q{Dot}$ e $\q{Err}$ já estão em blocos diferentes.
Como $\delta(\q{N},\texttt{.})=\q{Dot}$ e
$\delta(\q{F},\texttt{.})=\q{Err}$, o bloco $\{\q{N},\q{F}\}$ se divide:
\[
\mathcal{P}_2=\big\{\{\q{0}\},\{\q{I}\},\{\q{N}\},\{\q{Dot}\},
\{\q{F}\},\{\q{Op},\q{Sep}\},\{\q{W}\},\{\q{Err}\}\big\}.
\]
Todos os blocos agora são unitários, exceto $\{\q{Op},\q{Sep}\}$, cujos
destinos são sempre $\q{Err}$. Logo $\mathcal{P}_3=\mathcal{P}_2$.
O algoritmo termina com \textbf{oito estados}.

\subsection{Autômato quociente e preservação de categorias}
Seja $[q]$ o bloco que contém $q$ e $\q{T}=\{\q{Op},\q{Sep}\}$. Definimos
\[
M_{min}=(Q/{\sim},\Sigma,\delta_{min},[\q{0}],F_{min}),\quad
\delta_{min}([q],a)=[\delta(q,a)],
\]
onde $F_{min}=\{[\q{I}],[\q{N}],[\q{F}],\q{T},[\q{W}]\}$.
A transição está bem definida porque estados do mesmo bloco têm destinos
equivalentes. A tabela do quociente é obtida substituindo cada estado da tabela
original pelo seu bloco e unindo as linhas de $\q{Op}$ e $\q{Sep}$.
\begin{center}\begin{tikzpicture}[->,auto,node distance=17mm and 20mm,scale=.92,transform shape]
\node[state,initial] (s) {$[\q{0}]$};
\node[state,accepting,above right=of s] (i) {$[\q{I}]$};
\node[state,accepting,right=of s] (n) {$[\q{N}]$};
\node[state,right=of n] (d) {$[\q{Dot}]$};
\node[state,accepting,right=of d] (f) {$[\q{F}]$};
\node[state,accepting,below right=of s] (t) {$\q{T}$};
\node[state,accepting,below=of s] (w) {$[\q{W}]$};
\node[state,right=of t] (e) {$[\q{Err}]$};
\path (s) edge node {$L$} (i)
 (i) edge[loop above] node {$L,D$} (i)
 (s) edge node {$D$} (n)
 (n) edge[loop above] node {$D$} (n)
 (n) edge node {$.$} (d)
 (d) edge node {$D$} (f)
 (f) edge[loop above] node {$D$} (f)
 (s) edge node[right] {$O,S$} (t)
 (s) edge node[left] {$W$} (w)
 (w) edge[loop left] node {$W$} (w)
 (e) edge[loop right] node {$\Sigma$} (e);
\end{tikzpicture}
\end{center}
Também neste desenho, arestas omitidas levam a $[\q{Err}]$.

A fusão preserva a aceitação de lexemas. \textbf{Ela não preserva, sozinha, o rótulo
de saída OP/SEP.} Para classificar após a fusão, basta consultar o lexema:
se pertence a $O$, é OP; se pertence a $S$, é SEP. A implementação principal
mantém os nove estados para tornar a relação estado--categoria direta, e
calcula o quociente separadamente. Não afirmamos que um scanner com rótulos de
saída diferentes possa unir esses estados sem essa recuperação de categoria.

\subsection{Justificativa de minimalidade}
Entre estado final e não final, o sufixo $\varepsilon$ distingue o par. Entre os
três não finais, \texttt{a} distingue $[\q{0}]$ dos outros dois, e \texttt{0}
distingue $[\q{Dot}]$ de $[\q{Err}]$. Entre os cinco estados finais:
\begin{center}
\begin{tabular}{cl}
\toprule Pares & Sufixo que é aceito a partir de apenas um estado\\\midrule
$[\q{I}]$ contra qualquer outro final & \texttt{a}\\
$[\q{N}]$ contra $[\q{F}]$ & \texttt{.0}\\
$[\q{N}]$ ou $[\q{F}]$ contra $[\q{W}]$ ou $\q{T}$ & \texttt{0}\\
$[\q{W}]$ contra $\q{T}$ & espaço\\\bottomrule
\end{tabular}\end{center}
Todos os oito blocos são alcançáveis e dois a dois distinguíveis. Portanto,
nenhum outro par pode ser fundido preservando a linguagem: o AFD é mínimo.

\section{Do reconhecedor ao scanner}
O AFD recebe \emph{um lexema}; o scanner processa um arquivo completo.
Se o texto é \texttt{a+b}, o AFD rejeita a palavra inteira como um único lexema,
mas o scanner produz ID, OP e ID. Ele reinicia o AFD após cada token.
Um arquivo vazio produz uma lista vazia; isso não significa aceitar
$\varepsilon$ como um lexema.

\subsection{Maior prefixo aceito e retrocesso}
Em cada posição inicial, o scanner avança enquanto não chega ao estado de erro,
guardando a última posição em estado final. Ao parar, retorna a essa posição,
emite o token e reinicia. Se nenhuma posição final foi atingida, informa erro.
Essa é a política de maior casamento, habitual em análise léxica~\cite{aho}.

\begin{center}
\begin{tabular}{cll}
\toprule Entrada & Caminho relevante & Resultado\\\midrule
\texttt{12.50} & $\q{0}\to\q{N}\to\q{N}\to\q{Dot}\to\q{F}\to\q{F}$ & FLOAT\\
\texttt{ifx} & $\q{0}\to\q{I}\to\q{I}\to\q{I}$ & ID, após consulta\\
\texttt{12.} & último final antes do ponto & INT \texttt{12}; erro no ponto\\
\texttt{.5} & $\q{0}\to\q{Err}$ & erro na primeira posição\\\bottomrule
\end{tabular}\end{center}
Em \texttt{12.}, $\q{Dot}$ não é final: falta a parte fracionária. Após devolver
\texttt{12}, o scanner recomeça no ponto e acusa o erro. A função interrompe a
execução com exceção; não retorna a lista parcial de tokens.

\clearpage
\subsection{Trechos da implementação}
A tabela de transição do programa corresponde diretamente à especificação:
\lstinputlisting[firstline=23,lastline=33]{../../codigo/scanner.py}
O núcleo a seguir registra a última aceitação e executa o retrocesso:
\lstinputlisting[firstline=83,lastline=112]{../../codigo/scanner.py}
O código completo acompanha a entrega em \texttt{scanner.py}, sem bibliotecas
externas. A função \texttt{minimize()} implementa o refinamento mostrado no
relatório; \texttt{testes/test\_scanner.py} contém as verificações reproduzíveis.

\subsection{Revisão de software e rastreabilidade em 2026}
A versão 1.1.0 foi revisada em 13/11/2026. O código usa Python 3.12 ou superior,
anotações de tipos, aliases com a instrução \texttt{type}, registros imutáveis
com slots e uma interface de linha de comando com JSON. A documentação oficial
do Python 3.14 foi consultada na revisão~\cite{pythonatual}. A versão de Python
efetivamente testada está registrada em \texttt{evidencias/validacao.txt};
compatibilidade declarada não substitui teste em cada versão do interpretador.

Os recursos novos respondem a necessidades concretas: integração por JSON,
mensagens para arquivos inválidos, entrada padrão e preservação de quebras de
linha. Os conceitos de AFD, maior prefixo e Moore são clássicos e continuam
válidos. A data desta implementação não torna as técnicas novas; as datas
originais das referências são preservadas.

Na CLI, a tokenização completa ocorre antes da impressão: erros não produzem
uma lista parcial. \texttt{--json} apresenta tokens na saída padrão; erros
estruturados vão para a saída de erro. Os códigos de saída são 0 para sucesso,
1 para erro de arquivo/codificação/léxico e 2 para uso incorreto dos argumentos.
O arquivo padrão é resolvido relativamente ao código, não ao diretório corrente.

\subsection{Custo computacional e limites}
Neste AFD, há no máximo um caractere de prefixo consumido após o último estado
final: o ponto em $\q{Dot}$. O caractere que provoca a parada também pode ser
consultado novamente. Portanto o custo do scanner é $O(n)$, com quantidade
constante de consultas por caractere. Esse argumento é específico deste AFD;
não afirmamos esse limite para qualquer scanner com retrocesso.
O armazenamento dos tokens e lexemas é $O(n)$. As tabelas têm tamanho fixo.

O refinamento simples executado em Python possui, para $k$ classes e $m$
estados, no máximo $m$ rodadas e $O(km^2)$ avaliações de destinos no total,
desconsiderando custos de ordenação e estruturas auxiliares. Não se trata da
implementação otimizada de Hopcroft.

\section{Demonstração e verificação}
Execute os comandos na raiz do repositório, com Python 3.12 ou mais recente:
\begin{lstlisting}[language={}]
python codigo/scanner.py exemplos/exemplo.min --minimizar
python -m unittest discover -s testes -v
python codigo/scanner.py exemplos/invalido.min
\end{lstlisting}
O primeiro arquivo contém:
\lstinputlisting[language={}]{../../exemplos/exemplo.min}
Este é um exemplo para \emph{tokenização}; não se define nem se valida uma
gramática de comandos para esse texto.

Para a primeira linha, o resultado é:
\begin{center}
\begin{tabular}{clc}
\toprule Categoria & Lexema & Linha:coluna\\\midrule
KW\_LET & \texttt{let} & 1:1\\
ID & \texttt{total} & 1:5\\
OP & \texttt{=} & 1:11\\
FLOAT & \texttt{12.50} & 1:13\\
SEP & \texttt{;} & 1:18\\\bottomrule
\end{tabular}\end{center}
\texttt{invalido.min} apresenta erro no ponto de \texttt{12.}, em 1:15. O
encerramento com código 1 é esperado nesse teste negativo.

\subsection{O que os 18 testes verificam}
\begin{enumerate}
\item Categorias, maior casamento e operadores/separadores.
\item Todas as quatro palavras-chave.
\item Aceitação e rejeição de lexemas isolados, incluindo vazio e acentos.
\item Linha e coluna em múltiplas linhas.
\item Erros em \texttt{12.}, \texttt{.5}, \texttt{12.3.4} e \texttt{@}, com retrocesso.
\item Políticas explícitas para \texttt{12abc}, \texttt{-3}, \texttt{==}, espaços e arquivo vazio.
\item Totalidade da tabela e alcançabilidade dos nove estados.
\item Oito blocos, fusão correta e equivalência exata no produto alcançável.
\item Classificação de exatamente um caractere e comparação independente
com regex em 2.801 palavras representativas de comprimento até quatro.
\item Linha de comando: arquivo padrão independente do diretório, JSON,
entrada padrão, arquivo inexistente, UTF-8 inválido, erro sem saída parcial,
preservação de CR/LF e identificação da versão.
\end{enumerate}
Na verificação do produto, visitamos pares $(q,[q])$, comparamos aceitação e
seguimos as sete classes de transição. Como cada classe reúne caracteres com
o mesmo comportamento, ausência de divergência em todos os pares alcançáveis
prova equivalência para \emph{todas} as palavras, não apenas para exemplos de
comprimento limitado. Isso verifica o reconhecedor booleano; as categorias e a
política de tokenização são verificadas pelos testes específicos.

\section{Atendimento à proposta e conclusão}
\begin{center}\small
\begin{tabular}{p{.36\linewidth}p{.54\linewidth}}
\toprule Exigência & Evidência na entrega\\\midrule
Tupla matemática e transições & $M$, conjuntos e tabela completa na Seção 3\\
Diagrama limpo e legível & Original e mínimo, desenhados em TikZ\\
Minimização demonstrada & Alcançabilidade, $\mathcal{P}_0$ até $\mathcal{P}_3$, quociente e sufixos\\
Script Python/C++ e trechos & Python executável, trechos com listings e testes\\
Dois relatórios LaTeX & Técnico com article/TikZ/amsmath/listings; explicativo com foco no cliente\\
Slides LaTeX & Classe beamer; narrativa para aproximadamente 12 minutos\\
Defesa de 10--15 minutos & Apresentação planejada para aproximadamente 12 minutos\\\bottomrule
\end{tabular}\end{center}
A referência a árvores de derivação na rubrica é pertinente à alternativa com
gramáticas. Na aplicação escolhida, a evidência formal é a execução e a
minimização do AFD; não se apresenta uma árvore sintática como se fosse um
produto da análise léxica.

O projeto conecta expressões regulares, um AFD completo e um scanner executável.
A minimização reduz o reconhecedor de nove para oito estados e explicita a
diferença entre aceitar um lexema e atribuir-lhe uma categoria. As decisões
sobre fronteiras, decimais e erros tornam o comportamento reproduzível e
defensável oralmente.

\begin{thebibliography}{9}
\bibitem{sipser} SIPSER, Michael. \emph{Introduction to the Theory of Computation}.
3. ed. Cengage Learning, 2013.
\bibitem{aho} AHO, Alfred V.; LAM, Monica S.; SETHI, Ravi; ULLMAN, Jeffrey D.
\emph{Compilers: Principles, Techniques, and Tools}. 2. ed. Addison-Wesley, 2006.
\bibitem{pythonatual} PYTHON SOFTWARE FOUNDATION.
\emph{Python 3.14 Documentation: typing, dataclasses, argparse e unittest}.
Disponível em: \url{https://docs.python.org/3.14/}. Acesso em: 13 nov. 2026.
\end{thebibliography}
\end{document}
